\documentclass[12pt, a4paper]{report}
\usepackage[UTF8, heading=true]{ctex}
\usepackage{amsmath, amssymb, amsthm, mathrsfs}
\usepackage{pifont}
\usepackage{float}
\usepackage{geometry}
\usepackage{graphicx}
\usepackage{booktabs}
\usepackage[colorlinks=true, linkcolor=blue, citecolor=red]{hyperref}
% Prevent problems with linebreaks (\\) in PDF metadata strings
\pdfstringdefDisableCommands{\def\\{}}
\usepackage{hyperxmp}

\geometry{left=2.5cm, right=2.5cm, top=3cm, bottom=3cm}

\title{\textbf{大语言模型高级逻辑与语义缺陷研究：\\构造范式、根源剖析与评测基准扩展}}
\author{梁家瑞}
\date{2026年5月27日}

\begin{document}

\maketitle

\begin{abstract}

当前大语言模型在多模态理解与自然语言生成领域展现出令人惊叹的拟合表象，其在各类标准化推理测试中的卓越表现往往产生一种具备人类级别认知能力的错觉。然而，这种繁荣在本质上高度依赖于互联网巨量文本的模式匹配与浅层概率拟合，掩盖了底层心智模型深度的逻辑匮乏。

本研究聚焦于通过系统化、结构化、参数化的逆向工程，构造特定的逻辑与语义陷阱问题，系统测度并最终打破当前前沿大模型（如DeepSeek、ChatGPT、Claude）的概率预测边界。通过总结这些陷阱问题的特性，我们形成了一个 12 道题的 Benchmark，并测试了主流大模型在这上面的表现（见附录）。另外，本文还对模型犯错的原因进行了归纳，旨在深度透视 LLM 潜藏的认知缺陷，进而为下一代真正具备溯因推理体系的强人工智能指明可能的架构演进路径。
\end{abstract}

\tableofcontents
\newpage

\chapter{引言与大模型认知缺陷量化模型}

\section{研究背景及目标}

近年来，生成式预训练 Transformer 架构的爆发式演进，将人工智能推向了前所未有的视界。前沿大语言模型甚至在部分高难度学科竞赛、司法考试以及医疗基准测试中逼近或超越了人类专家的平均基准。这种在标准评测体系下的繁荣，促使大量研究者乐观地预言通用人工智能的奇点即将来临。然而，这种基于庞大参数量级的成功，其底层逻辑依托的极其纯粹——海量语料参数化映射后的高维向量内积与自回归统计学预测。大语言模型当前展现的诸多“智慧”表征，脱胎于海量文本训练形成的参数权重，本质上属于极致的经验记忆与模式复刻。模型在此展现的绝未触及思维本质，其内核仅仅停留在高维向量空间的概率映射。

当测试用例稍微偏离主流经验分布的流形，进入长尾逻辑空间时，这些庞大的算力巨兽往往会暴露出极其荒谬的逻辑断裂现象。本研究通过大量的实证分析与逆向探索发现，大模型的诸多低级错误背后，潜藏着高度一致的运作机制缺陷。这种缺陷源于其对长上下文注意力分配的失衡，以及面对图论、数论等隐性属性校验时结构性溯因能力的彻底缺失。

本研究的根本核心定位于通过系统化、参数化的方法构造特定结构的逻辑/语义陷阱问题。此类问题犹如针对神经网络架构的探针，精密测试并试图打破大模型的概率预测边界。在这个过程中，我们将抽丝剥茧地总结出模型在面对多约束冲突或资源枯竭时犯错的共性特征。更进一步，我们寄望于通过透视当前LLM的心智模型缺陷，针对诸如注意力矩阵坍塌、概念重叠解耦失败等现象，为下一代寄望具备真正非凡认知与严密演绎能力的AI系统，梳理出一条迥异于纯粹规模扩张的进化路径。

\section{逻辑困局的核心量化评测体系}

欲使对AI认知缺陷的诊断上升至严谨的学术层级并具备普适的指导原则，我们必须摒弃定性的模糊描述，转而将大语言模型的犯错机理提炼为可复现、可操控的核心认知参数。基于此，本章引入两个至关重要的量化指标，用以严格界定陷阱问题构造的理论难度，并精准测度大模型陷入逻辑谬误的致错概率。

\subsection{条件先验概率依赖度（Conditional Prior Probability Dependency, CP）}

在大模型的生成范式中，Token的输出始终受到自回归条件分布的严密支配。定义某一高阶问题 $B$（例如一个无向图是否有欧拉回路），和在讨论这个问题的语境下需要面临的低阶问题 $A$（例如这个无向图是否联通）。我们将模型对于高阶问题逻辑链路推演的敏感性，具象化为条件先验概率依赖度 $CP= P(A|B)$。

当该条件概率极度逼近必然，即 $CP \approx 1$ 时（例如在数学教科书以及海量算法竞赛真题库中，几乎所有讨论无向图是否有欧拉回路的问题都着重探讨顶点度数的奇偶性，反而图的联通性是显然的），大模型在进行推理演算时将产生强烈的概率路径依赖。这种依赖机制宛如一种思维毒药。在计算图论节点的度数消耗了模型大量局部算力与上下文窗口后，系统会面临极高的置信度诱惑，进而直接跨越对隐性核心属性（如图连通性）的严谨拓扑验证，鲁莽且毫无根据地给出一个形式上看似无懈可击逻辑上却毫无道理的谬误结论。

这种基于统计频率对齐的“条件反射”，生动折射出LLM缺乏多步验证机制的心智缺陷。在人类数学家的认知框架内，必要条件的满足必定激发对充分条件严阵以待的审查；而在参数化的大模型视域下，高频共现的Token序列宛如深不见底的引力井，轻易将尚处于演化状态中的逻辑链条强行揉捏成顺滑却空洞的断言。

\subsection{算力/洞察非对称指数（Computational vs Insight Disparity, CD）}

大模型在处理深层数理结构时面临的第二重困境，源于其离散推理步骤与其有限计算资源的不可调和性。为了精确量化这种计算隔阂，我们设计了算力/洞察非对称指数（CD指数）。

设有特定数学验证题的某个隐藏属性，欲通过穷举模拟或基础启发式搜索证明该属性（例如借助深度优先搜索 DFS 或广度优先搜索 BFS 在高维线性空间或庞大 Cayley 图中暴力遍历检查网络连通性），所需的基础计算复杂度定义为 $C_{bf}$。同时，若引入某种具备深层结构美感的数学对称性、组合守恒量或近世代数定理直接破解该问题，其“洞察复杂度”定义为 $C_{in}$。由此，非对称指数 $CD$ 获得如下严格定义：
\begin{equation}
CD = \log \left( \frac{C_{bf}}{C_{in}} \right)
\end{equation}

对于一个阈值 $\text{Threshold}$，当构造的问题满足 $CD > \text{Threshold}$ 且 $C_{bf}$ 在物理层级彻底突破了大语言模型的单次推理 Token 缓冲区极限或最大时间步长限度时，该命题对模型而言即进入了广义的非多项式（NP）验证深渊。反观人类数学直觉或纯粹代数视角下，借助几何对称性或不变量理论等“观察”，其 $C_{in}$ 仅仅为 $O(1)$ 级别。面临此类极端非对称性博弈时，模型系统由于无法承担起指数爆炸般的暴力演算，内部潜藏的正反馈奖励对齐倾向又剥夺了其宣告“无解”或“超出算力极限”的勇气。最终，模型缺乏提取理论守恒量所需的真正数学洞察机制，被硬性驱逐进自我麻醉的幻觉生成状态，选择性且极具欺骗性地“口糊”跳过核心逻辑节点的校验程序，编造出荒诞不经的逻辑链完备性假象。

\section{本章小结}

通过上述的 $CP$ 指标与 $CD$ 指数，我们已经初步完成完成对LLM心智缺陷侦测体系在理论参数维度的框架搭建。这两个量化抓手揭开了模型看似精密逻辑外衣下高度依赖经验匹配的虚弱本质。下一章我们将基于该理论体系展开，通过若干个极具破坏美学的陷阱命题构建（涉及从基础语义状态纠缠到高阶近世代数超约束拓扑等多个维度），解剖大模型如何在自身注意力矩阵衰变与高频语境裹挟的双重倾轧下，走向逻辑验证的崩溃。

\chapter{认知陷阱的构造范式与底层机理剖析}

在明确了模型缺陷的量化体系之后，本章将实地铺设“陷阱网络”。我们将以多维度的跨界语汇、长程链条以及高维图论网络为试验道场，全面演示大模型智能表象下的认知塌缩过程。我们不仅复盘错误的演化时刻，更着力提炼可供无限复制的 Benchmark 构造范式。

\section{语义匹配与惯性模式识别陷阱：隐形债务与无形现金流错位}

金融交易与日常借贷场景历来是测试物理世界抽象能力的绝佳切口。人类在面对此类问题时，往往通过内建的心智状态机，将资产与债务在独立的逻辑账本中隔离处理。大语言模型处理此类文本时，由于缺乏强制状态多维解耦机制，频繁在多实体间沦陷。

\subsection{基准测试的原题重现与深度拆解}

\textbf{[问题 1.0]} \\
小明去理发店理发，理发需要70元。可是他搞忘带钱了，于是他向理发师借了140元，并将这140元中的70元支付给了理发师当作理发费用。请问他还欠理发师多少元？

\textbf{[标准答案]}：140元。

\textbf{[标准数学/逻辑推导]} \\
该问题的本质在于区分“服务类无形消耗所生成的隐性负债”与“实体货币现金流注资所生成的显性债务”。我们可通过分步财务状态方程进行严谨推演：
\begin{itemize}
    \item \textbf{时刻 0（初始基线）}：小明的总资产 $A_0=0$，总负债 $L_0=0$；理发师对应的债权 $C_0=0$。
    \item \textbf{时刻 1（接受服务，产生隐性负债）}：小明完整消耗了标价为 70 元的理发服务，此举即刻转化为商业契约下的硬性债务。此时小明负债 $L_1 = +70$。
    \item \textbf{时刻 2（现金流注入）}：发生 140 元的借贷行为。小明此时现金增广 140 元，然而其总债务须叠加此笔纯粹的借贷本金。总债务更新为 $L_2 = 70 + 140 = 210$ 元。
    \item \textbf{时刻 3（支付冲抵）}：小明流出 70 元现金交付理发师，专项旨在平抑因理发服务遗留的债务账单。此时总负债削减 70 元。故最终的稳态债务为 $L_3 = 210 - 70 = 140$ 元。
\end{itemize}
于是答案是140元。

\subsection{大模型的回答与微观致错机理剖析}

在前沿大模型（如 DeepSeek、Claude Sonnet）的首轮零样本测试中，系统展现出了令人咋舌的脆弱性，高频次且极其笃定地输出“70元”这一谬见。当迫使模型分步写出思维链时，系统的遮羞布方被彻底扯下。

探其深层认知科学病理，LLM本质上正经历一场严重的局部语义黑洞吞噬现象。受高维语料库的惯性训练裹挟，“支付”、“借款”、“结欠”等强关联金融词汇构建了一个极易落入的浅层引力场。在人类普遍语流中，“借去 $2X$，随之支付 $X$”这句话后续最为常见的概率接续文本，几乎毫无悬念地指向“还欠 $X$”。
模型被这种粗暴的模式流劫持，最终导致多实体归属混淆以及隐性负债的选择性无视。它无法将“理发无形服务开销”（实质负债）与“物理现金流向”（账面结款）在抽象符号层面上进行正交解耦，注意力矩阵坍塌至局部近义词的强行加减运算中。

\subsection{构造范式的参数化抽象与批量拓展}

这一发现为系统性挖掘AI盲区提供了强有力的武器。我们可以通过如下几套固定约束机制，参数化生成海量Benchmark：
\begin{enumerate}
    \item \textbf{刻意植入隐性消耗}：虚构一类日常行为，含有不易立刻激发“借据”敏感性的隐式耗损（量化为 $X$）。
    \item \textbf{施加现金流掩护与双倍重叠}：设计债权方反向注资 $2X$ 的物理货币，并促使受资方原封不动地用此笔款项拆解覆盖先前的无形支出。
    \item \textbf{密集轰炸语义干扰源}：高密度拼接“递给”、“找零”、“顺手归还”等强情绪或者动态画面扰动词，迫使模型在计算最后一步时，Token 序列急速向 $\text{借入 } 2X - \text{交付 } X = X$ 的残缺代数系统坠落。
\end{enumerate}

这种范式精密对应了上文提及的条件概率诱导机制。以下将进一步展示几道依照该参数模板扩展所得的标准 Benchmark 测试题。

\textbf{[问题 1.1 酒店入住与押金错位问题]}：小张入住一家商务酒店，房费每晚300元。由于出门匆忙，小张忘带钱包和手机。恰好他在大堂遇到了熟人李总，于是小张向李总借了600元现金。随后，小张将这600元现金全部交给了酒店前台，其中300元用于支付当晚房费，另外300元作为入住押金。第二天中午退房时，前台检查房间无损，将300元押金以现金形式退还给了小张。小张在大堂刚好又遇到了李总，便顺手把刚刚退回的300元现金还给了李总。请问此时小张还欠李总多少钱？

\textbf{[标准答案]}：300元。

\textbf{[标准推导与致错机理]}：小张的核心负债总计为 600 元。最后退还给李总的 300 元仅抵偿了借款总额的一半，另外 300 元早已作为不可退还的房费被消耗。因此，小张仍欠李总 300 元。大模型在处理该测试用例时，其注意力极易发生局部聚类，被“酒店完整退还押金，小张顺手归还友人”这一表面高度自洽的闭环语义所误导。模型倾向于生成“小张借入600，退回300，再还300，债务彻底平息”的谬误输出，彻底遗漏了最初借款中已有 300 元已转化为沉没状态的服务成本。

\textbf{[问题 1.2 餐馆 AA 制与替付罗生门问题]}：甲、乙、丙三人去餐馆聚餐，结账时总共花费300元，三人约定完全AA制，每人应承担100元。然而甲结账时发现自己手机没电无法支付，乙便主动提出替甲垫付这100元。因此，在收银台前，乙一共通过微信支付了200元，丙支付了100元。聚餐结束离开餐馆后，甲觉得过意不去，便向丙借了100元现金，并当场将这100元现金递给了乙，表示把刚才饭钱还上。请问此时，甲、乙、丙三人之间的最终债务清偿关系是怎样的？谁欠谁多少钱？

\textbf{[标准答案]}：甲欠丙100元；甲与乙之间、乙与丙之间均不存在任何债务关系。

\textbf{[标准推导与致错机理]}：本题的债务网络非常清晰。甲应对聚餐承担 100 元，借助乙的垫付完成结账；随后，甲向丙借款 100 元还清了乙的垫资。最终的债务关系为：甲欠丙 100 元。甲与乙、乙与丙之间均不存在任何债务关系。然而，在处理此类问题时，题目中穿插的“垫付”、“借款”、“还钱”等多向实体网络会引发严重的解耦失败。大模型在追踪“甲 $\rightarrow$ 乙（垫付）”、“丙 $\rightarrow$ 甲（借款）”、“甲 $\rightarrow$ 乙（还钱）”等动态拓扑时，无法将各实体的账本独立开来，导致注意力机制紊乱，常常得出“乙还欠丙钱”或“甲依然欠乙100元”的错误结论。

\textbf{[问题 1.3 二手车交易与定金套现陷阱]}：老王看中了小李的一辆价值2万元的二手摩托车。老王手头没有现金，小李慷慨地借给老王4万元现金。老王接过这4万元现金后，立刻将其中的2万元作为车款付给小李，把车骑走；同时，老王为了表示诚意，将剩下的2万元现金也交还给小李，作为该交易的长期履约定金（合同约定定金可随时全额退还）。三天后，老王因家人反对决定退车。根据合同，小李同意退车，并当场把2万元履约定金全额退还给了老王。老王拿到这2万元现金后，心里觉得很不好意思，立刻把这2万元现金塞回给小李，说：“这2万元先还给你当借款。” 随后老王把摩托车完好无损地推回给小李。请问此时老王还欠小李多少钱？

\textbf{[标准答案]}：0元。

\textbf{[标准推导与致错机理]}：该命题构造了高阶的现金流双向重叠。起初，老王的借款总额为4万元。退车消除了2万元的车款资产，定金的退回消除了2万元的定金资产，故最终两清。大模型在面对“摩托车推回、定金退还并塞回”的连续高情绪扰动词汇时，会产生思维惰性，无法分清哪一步是真正的债务关系哪一步只是单纯的购买。


\section{逻辑长程推进后的失忆与逃脱：超约束代数体系下的隐秘崩塌}

大模型展现出惊人的一维片段提取能力，但在长距离、多阶段、强逻辑耦合之下的连续推导过程中，其强化学习惩罚驱动架构极度容易促发系统的主动逃避或历史条件的蓄意遗忘。

\subsection{基准测试的原题重现与深度拆解}

\textbf{[问题 2.0]} \\
设 $A$ 是实数域 $\mathbb{R}$ 上的一个 $4 \times 4$ 矩阵。已知 $A$ 满足以下条件：
\begin{enumerate}
    \item 矩阵 $A - 2I$ 的零空间维度为 1。
    \item 矩阵 $(A - 2I)^2$ 的零空间维度为 2。
    \item 存在一个 2 维的不变子空间 $V \subset \mathbb{R}^4$，使得对于任意向量 $v \in V$，都有 $(A^2 + A + I)v = 0$。
    \item 已知向量 $v_1, v_2, v_3$ 线性无关，且满足 $(A-2I)v_1 = 0$, $(A-2I)v_2 = v_1$, $(A-2I)v_3 = v_2$。
\end{enumerate}
请根据以上条件，严格推导并写出矩阵 $A$ 在复数域 $\mathbb{C}$ 上的若尔当标准型。

\textbf{[标准答案]}：此题无解（条件自相矛盾，满足该条件的 $4 \times 4$ 矩阵根本不存在）。

\textbf{[标准长程推演与断裂暴露]} \\
依循经典线性代数原则，此题实质上被封装于超约束的系统中：
\begin{itemize}
    \item \textbf{链条前端推导}：由条件 4 可知，关于特征值 $\lambda = 2$ 存在一条长度至少为 3 的若尔当链。这硬性界定了 $\lambda = 2$ 对应的若尔当块的规模至少为 $3 \times 3$。结合条件 1（几何重数受限为 1），可以确定关于 $\lambda = 2$ 的若尔当块独占了系统中的 3 到 4 个维度。
    \item \textbf{中端扰动剖析}：条件 3 引入了在二维不变子空间 $V$ 中的多项式方程 $x^2 + x + 1 = 0$。它在复数域中必然存在一对共轭复根 $\omega, \bar{\omega} = \frac{-1 \pm \sqrt{3}i}{2}$。由于基础矩阵为实矩阵，为容纳这对共轭复根，必须分配 2 个独立的 $1 \times 1$ 复若尔当块。
    \item \textbf{维度极限挤压}：该矩阵复化后所需的总空间维度已达：关于 $\lambda = 2$ 的若尔当块需求（$\ge 3$）加上共轭复特征值占据的需求（$2$），总计 $\ge 5$ 维。这与系统基设前提—— $A$ 为一个 $4 \times 4$ 矩阵，发生了不可调和的空间冲突。因此，满足条件的 $4 \times 4$ 矩阵根本不存在。
\end{itemize}

\subsection{大模型的回答与微观致错机理剖析}

当 Deepseek 在处理此类包含冲突陷阱的命题时，我们观察到了其核心认知缺陷：长距逻辑推理下的多约束一致性丧失与“逃避心理”。

在推理初期，大模型能够精准地将条件 3 翻译为“包含 2 个 $1 \times 1$ 的共轭复块”，也能将条件 4 翻译为“包含一个至少 $3 \times 3$ 的若尔当块”。然而，当局部推演得出的总维度需求（$3+2=5$）与全局物理前提（$4 \times 4$ 矩阵）发生不可调和的冲突时，大模型并未展现出人类的“溯因推理”反向校验能力。在强化学习的对齐机制下，模型往往将“输出无解”视为可能招致惩罚的负收益行为。

为了迎合并强制输出一个形式上完整的答案，模型在经过大量计算之后，选择了“遗忘”或主动抛弃条件 4 中关于若尔当链长度的限制，极具欺骗性地强行拼凑出一个包含 $2 \times 2$ 若尔当块及共轭复特征值的 $4 \times 4$ 矩阵。这种现象本质上暴露了大模型在长上下文状态下无法保持严密逻辑自洽的深层缺陷。

\subsection{构造范式的参数化抽象与批量拓展}

这种深层缺陷为我们提供了系统性构造测试用例的抽象范式。我们可以通过以下规则参数化生成高阶超约束 Benchmark 测试题：
\begin{enumerate}
    \item \textbf{构建超约束系统}：设计一道涵盖线性代数、近世代数或微积分的高阶理论推导题，需经过 4 步以上的演绎。
    \item \textbf{多源维度挤压}：通过前置条件 A、B、C 定义局部硬性性质，隐式地将代数系统或几何空间的维度（阶数）上限压缩至 $N$。
    \item \textbf{长链末端注入冲突}：在题干末尾，刻意植入一个形式上十分自然、独立看似合理的条件 D。然而该条件本身隐含的理论下限要求至少为 $N+1$ 空间。由此导致整个多约束代数系统彻底崩溃并绝对无解。
\end{enumerate}

\textbf{[问题 2.1 三维向量空间与隐形线性无关冲突]}：设 $V$ 是实数域 $\mathbb{R}$ 上的一个 3 维线性空间，$\alpha_1, \alpha_2, \alpha_3, \alpha_4$ 是 $V$ 中的四个非零向量。已知它们满足以下四个条件：
\begin{enumerate}
    \item 向量组 $\alpha_1, \alpha_2, \alpha_3$ 严格线性无关；
    \item 向量 $\alpha_4$ 无法由向量组 $\alpha_1, \alpha_2$ 线性表出；
    \item 存在 $V$ 的一个 2 维子空间 $W \subset V$，使得 $\alpha_3 \in W$ 且 $\alpha_4 \in W$；
    \item 存在 $V$ 上的一个线性变换 $T$，使得该变换在四个向量上满足循环映射关系：$T(\alpha_1) = \alpha_2, T(\alpha_2) = \alpha_3, T(\alpha_3) = \alpha_4, T(\alpha_4) = \alpha_1$。
\end{enumerate}
请严格推导并求出线性变换 $T$ 在基 $\{\alpha_1, \alpha_2, \alpha_3\}$ 下的矩阵表示。

\textbf{[标准答案与推导]}：此题无解。条件4暗示着在这四个向量组张成的子空间上 $(T^4 - 1) = (T-1)(T+1)(T+i)(T-i) = 0$，这四个向量如果想在循环映射下满足所有条件而不退化，它们在空间中对应的几何轨迹是一个完整的 4 阶循环，要张成一个能容纳这个 4 阶循环映射的最小空间，其特征多项式必须包含 $(T^4-1)$ 的全部因子，于是一个三维的线性变换拥有了四个不同的特征值，矛盾。

\textbf{[致错机理]}：AI在面对此题时，会投入大量Token去建立变换矩阵，并利用条件1, 2, 3去解齐次线性方程组。在推导超过1000个Token后，由于上下文的注意力衰减，它会忽略基向量数量（3个）与映射链条（4个向量独立）之间的空间维数根本冲突，最终强行杜撰出一个 $3 \times 3$ 的错误矩阵。

\textbf{[问题 2.2 有限阿贝尔群阶数与子群结构的隐式冲突]}：设 $G$ 是一个有限阿贝尔群（Abelian Group）。已知该群同时满足以下四个性质：
\begin{enumerate}
    \item 群 $G$ 的阶数严格为 $|G| = 12$；
    \item 群 $G$ 中存在一个阶数为 4 的循环子群 $H \cong \mathbb{Z}_4$；
    \item 群 $G$ 中不包含任何阶数为 6 的元素；
    \item 群 $G$ 包含一个同构于克莱因四元群（Klein 4-group）的子群 $K \cong \mathbb{Z}_2 \times \mathbb{Z}_2$。
\end{enumerate}
请根据阿贝尔群有限生成基本定理，写出群 $G$ 的同构分类及直积分解形式。

\textbf{[标准答案]}：满足条件的有限群不存在，题目条件自相矛盾。

\textbf{[标准推导与致错机理]}：根据有限阿贝尔群分类定理，12阶群只有两种：$\mathbb{Z}_{12}$ 和 $\mathbb{Z}_2 \times \mathbb{Z}_6$。若 $G = \mathbb{Z}_{12}$，则包含4阶循环子群（由3生成的子群），但其本身也包含6阶元素（由2生成的子群），与条件3冲突。若 $G = \mathbb{Z}_2 \times \mathbb{Z}_6$，其中元素的最高阶为6（与条件3冲突），且该群中所有元素的阶均无法达到4，故不含4阶循环子群。大模型在处理这类多约束代数题时，往往会在前几步成功列出12阶群的分类，但在长距离推理中，会由于无法同时满足 $\mathbb{Z}_4$ 与不含6阶元素的约束，最终选择主动忽略条件3，错误地给出 $G \cong \mathbb{Z}_{12}$ 的答案。

\textbf{[问题 2.3 欧氏空间正交变换与迹的隐式绝杀]}：设 $V$ 是 3 维实欧几里得空间，$T$ 是 $V$ 上的一个正交变换，且该变换的行列式满足 $\det(T) = 1$。已知 $T$ 还满足下列性质：
\begin{enumerate}
    \item 变换 $T$ 的迹（矩阵所有对角元素之和）满足 $\text{Tr}(T) = -1$；
    \item 存在 $V$ 的一个 2 维不变子空间 $W \subset V$，使得 $T$ 限制在 $W$ 上的诱导变换 $T|_W: W \rightarrow W$ 是一个旋转角为 $\pi/3$ 的平面旋转；
    \item 变换 $T$ 不是恒等变换。
\end{enumerate}
请严格求出变换 $T$ 的特征多项式，并写出其在某组标准正交基下的矩阵。

\textbf{[标准答案]}：条件自相矛盾，此类正交变换在物理和数学上均不存在。

\textbf{[标准推导与致错机理]}：由于 $T$ 是 3 维正交变换且 $\det(T)=1$，根据欧拉旋转定理，$T$ 必有一个特征值为 1（对应旋转轴）。另外两个特征值来自平面旋转 $W$，其旋转角为 $\pi/3$，故对应的复特征值为 $e^{\pm i\pi/3} = \cos(\pi/3) \pm i\sin(\pi/3) = \frac{1}{2} \pm \frac{\sqrt{3}}{2}i$。因此，该正交变换在复数域上的全部三个特征值为 $\{1, e^{i\pi/3}, e^{-i\pi/3}\}$。其迹 $\text{Tr}(T)$ 必须等于所有特征值之和，即 $1 + 2\cos(\pi/3) = 1 + 1 = 2$。这与条件1规定的 $\text{Tr}(T) = -1$ 发生了彻底的几何冲突。大模型会深陷于写出三维旋转矩阵的形式表达，并在漫长的三角函数几何代数推导中，将最初的“迹为 -1”这一前置条件抛诸脑后，给出基于旋转角 $\pi/3$ 的错误矩阵。


\section{高先验概率陷阱与算力不对称：以 Cayley 图的隐性不连通现象为例}

在处理图论或离散结构验证问题时，大模型高度依赖基于过往频率分布的启发式结论。当遇到计算复杂度呈现指数级跃升，而理论校验仅需特定代数洞察即可化解的任务时，模型陷入逻辑灾难的过程尤为典型。此类缺陷集中暴露了当前 LLM 尚未跨越算力壁垒与深层代数恒等量理解之间的认知鸿沟。

\subsection{基准测试的原题重现与深度拆解}

\textbf{[问题 3.0]} \\
考虑一个拥有 $2^8 = 256$ 个顶点的无向图 $G$，顶点由所有的 8 位二进制向量 $v \in \{0, 1\}^8$ 表示。定义该图的边集生成规则如下：两个顶点 $u, v$ 之间存在一条边，当且仅当 $u \oplus v \in S$（其中 $\oplus$ 为按位异或运算）。生成集 $S$ 定义为：
$$S = \{ x \in \{0, 1\}^8 \mid \text{popcount}(x) \equiv 0 \pmod 2 \text{ 且 } x \text{ 中包含且仅包含两个连续的 1 代码块} \}$$
（注：例如 01100110 属于 $S$，因为它包含两块连续的 1，且 popcount 为 4；而 11100000 不属于 $S$。）
问题：请推导并判断该图 $G$ 是否存在欧拉回路（Eulerian Circuit）？

\textbf{[标准答案]}：不存在欧拉回路（因为该图 $G$ 是不连通的，存在多个完全隔离的孤立连通分支）。

\textbf{[标准数学推导]} \\
此问题需经历严格的双阶段拓扑验证。
\begin{itemize}
    \item \textbf{第一阶段（显性算力消耗度数校验）}：在此特定 Cayley 图中，每一个顶点的度数均精确等效于生成集 $S$ 的基数大小（即 $\text{deg}(v) = |S|$）。通过完备的组合计数分析，筛选出恰好包含两块 1 且包含偶数个 1 的向量，可测算得出 $|S| = 66$。由于 66 为偶数，全图顶点的度数奇偶性条件完美达成。
    \item \textbf{第二阶段（隐性不变量与连通性洞察）}：存在欧拉回路要求全图必须连通。根据定义集 $S$ 的内在刚性约束，所有合法映射模式均服从 $\text{popcount}(x) \equiv 0 \pmod 2$ 标准。因而，从任何一顶点 $u$ 循边游走至相邻顶点 $v = u \oplus s$，新顶点 $v$ 中 1 的比特数将保持与 $u$ 一致的奇偶性。这一底层同余守恒律将整个 256 节点的布尔空间劈裂为两座“平行的孤岛”。具有偶数个 1 的顶点集与拥有奇数个 1 的顶点集之间无法互相到达。图 $G$ 因此不连通而不存在欧拉回路。
\end{itemize}

\subsection{大模型的回答与微观致错机理剖析}

在应对此类陷阱时，ChatGPT 回答存在欧拉回路。其运算分配的轨迹彻底暴露出算力/洞察非对称指数（CD）失衡引发的灾难。模型将超乎寻常的绝大多数 Token 投入到第一阶段 $|S|$ 集合元素的强行罗列与宏大组合数学推理中。但在这场巨大的算力消耗落幕后，它对第二阶段的连通性校验极其草率，悍然输出诸如“因 $S$ 基数庞大且形态丰富，自然具备连通整个空间的生成张力，故该图连通……”的荒谬引理。

这种崩溃源自行文之初所述的高先验概率陷阱。在训练语料库中，$CP \approx 1$ 导致条件依存度趋近于绝对存在。这种深刻的统计印记诱发了致命的概率顺从性。加以验证图是否真正连通本身存在着极其高昂的基础算法复杂度（DFS/BFS爆搜时间复杂度 $O(4^n)$ ），模型无法在单次序列生成过程中承载。在缺乏寻找 $O(1)$ 常数级“奇偶不变量”等高级代数直觉的心智空白地带，为了消解算力真空期带来的注意力断裂与概率焦段失焦问题，大模型极其自发地选用了最具备文字连续流畅度的经验结论以瞒天过海。

\subsection{构造范式的参数化抽象与批量拓展}

以此类算力与洞察不对称的陷阱为蓝本，我们总结出如下极具穿透力的参量化构造法则：
\begin{enumerate}
    \item \textbf{设计双阶段验证失衡题}：迫使系统首先进入一个高度显性、极度消耗推演步骤的代数或组合计算区间；随后埋伏一个牵涉拓扑连通、空间非空或全局同构性的验证环节。
    \item \textbf{隐性守恒量植入}：在看似随机或复杂的生成集与状态转移矩阵中，不动声色地注入一个底层的代数守恒律（如特定的模 $N$ 同余、比特位锁定、理想包含律等）。该守恒律将在暗中将庞大的状态空间强行切割为若干不可逾越的孤立轨道。
    \item \textbf{调用高先验概率施行麻痹}：刻意操纵第一阶段耗费算力所得出的结果完美契合某项著名数学定理的前置范式，从而引发模型极高的概率对齐惯性，以此彻底麻痹其对第二阶段非平凡属性的独立审查机制。
\end{enumerate}

\textbf{[问题 3.1 棋盘翻转奇偶性不一致的陷阱]}：在一个无限大的棋盘上，一个格子里放置着一个标准正方体。每次操作可以控制该正方体向前、后、左、右四个相邻的格子之一翻滚 $90^\circ$（即绕正方体底部的某条棱滚动到相邻格子上）。为了精确描述正方体的姿态，我们将一面印上奶龙的脸（有明确的“上下”方向，以其嘴巴所在位置为基准）。初始奶龙朝向前方，嘴巴向下。请问，如何操作可以使得正方体最终回到原来的格子，并且奶龙朝向上方，嘴巴朝向前方？

\textbf{[标准答案]}：无解。

\textbf{[标准推导与致错机理]}：一方面，最终回到原点意味着经历了偶数次翻转。另一方面，考虑对正方体的六个面各选择一条对角线连接，并且这六条线构成一个正四面体。从上方俯视正方体时，顶面的对角线要么从左下到右上（记为偶状态），要么从左上到右下（记为奇状态）。那么每进行一次翻转，正方体状态的奇偶性都会发生变化。注意到初始状态和目标状态的奇偶性不一致，故需要奇数次翻转。综上所述，本题无解。本题的关键在于棋盘无限大，强行验证是理论上不可行的。于是，为了迎合用户的需求，大模型会毫无根据地胡乱描述一通并最终给出一套完全错误的方案。

\textbf{[问题 3.2 高阶隐式同余下平稳分布的绝杀陷阱]}：设一个具有 100 个状态的齐次马尔可夫链，状态空间由整数集合 $S = \{1, 2, \dots, 100\}$ 表示。对于状态空间中的任意状态 $i \in S$，我们将其唯一定位到一个 $10 \times 10$ 的二维坐标上：令 $i - 1 = 10x_i + y_i$，其中 $x_i \in \{0, 1, \dots, 9\}$ 代表其行坐标，$y_i \in \{0, 1, \dots, 9\}$ 代表其列坐标。

该马尔可夫链的一步转移概率矩阵 $P = (p_{ij})_{100 \times 100}$ 的构造规则如下：
从状态 $i(x_i, y_i)$ 转移到状态 $j(x_j, y_j)$，设行跨度为 $\Delta x = (x_j - x_i) \pmod{10}$，列跨度为 $\Delta y = (y_j - y_i) \pmod{10}$。
\begin{itemize}
    \item 若 $\Delta x = 0$ 且 $\Delta y = 0$（即状态试图转移到自身），则转移概率严格为 $p_{ii} = 0$；
    \item 若 $(\Delta x, \Delta y) \neq (0,0)$，且满足以下三次齐次同余方程：
    $$\Delta x^3 + 4\Delta x^2\Delta y + 2\Delta x\Delta y^2 + 2\Delta y^3 \equiv 0 \pmod 5$$
    则状态 $i$ 可以向状态 $j$ 转移，且对于状态 $i$ 出发的所有合法目标状态 $j$，其转移概率 $p_{ij}$ 严格相等。
    \item 若不满足上述方程，则转移概率 $p_{ij} = 0$。
\end{itemize}
  1. 请严格证明该转移概率矩阵 $P$ 是否为双斯托卡斯特矩阵（Double Stochastic Matrix，即每一行和每一列的元素求和均严格等于 1）；\\
  2. 请推导并判断该马尔可夫链是否存在唯一的平稳分布（Stationary Distribution），且无论初始状态如何，长时间运行后该链都会收敛到该分布？

\textbf{[标准答案]}：1. 是双随机矩阵。2. 不存在唯一的平稳分布，也无法从任意初始状态收敛全局（整个状态空间被锁死分裂为 5 个平行的孤立子空间）。

\textbf{[标准推导与致错机理]}：证明其为双随机矩阵需利用阿贝尔群 $\mathbb{Z}_{10} \times \mathbb{Z}_{10}$ 上的平移不变性：无论正向扩散或逆向溯源，合法位移解集在全空间内高度几何对称，故每行与每列的最终求和必然等价为 1。然而，系统的连通性隐藏着绝杀陷阱。对复杂的三次同余方程在有限域 $\mathbb{F}_5$ 上进行因式分解，可得 $\Delta x^3 + 4\Delta x^2\Delta y + 2\Delta x\Delta y^2 + 2\Delta y^3 \equiv (\Delta x - 2\Delta y)^3 \equiv 0 \pmod 5$。这等价于线性方程 $\Delta x \equiv 2\Delta y \pmod 5$，同乘 3 移项即为 $\Delta x + 3\Delta y \equiv 0 \pmod 5$。这意味着转移发生时状态算子变化规律必须服从：$x_j + 3y_j \equiv x_i + 3y_i \pmod 5$。因此，在任何游走中，代数特征值 $C(i) = (x_i + 3y_i) \pmod 5$ 成为被严格锁死的隐性不变量。全局 100 个状态被冷酷地切割为 5 枚分别包含 20 个元素的拓扑孤岛。该离散网络根本丧失了不可约性。大模型面对双随机网络问题时极易陷入高级习题常见的“状态遍历与唯一定理”统计先验中。受限于非对称算力壁垒，大模型缺乏跳出文字泥沼进行抽象多项式分解的心智，便主观顺应文字诱导断言“极其丰富的非线性转移域必定使得空间连通”，在毫无觉察下惨遭理论滑铁卢。

\textbf{[问题 3.3 近世代数——满同态的隐形空集陷阱]}：设 $\mathbb{Z}[\sqrt{3}] = \{ a + b\sqrt{3} \mid a, b \in \mathbb{Z} \}$ 是实数域上的一个子环，令 $U$ 为该环的单位群（乘法可逆元群）。我们知道，一个元素 $x = a + b\sqrt{3} \in \mathbb{Z}[\sqrt{3}]$ 属于 $U$ 当且仅当其代数范数 $N(x) = a^2 - 3b^2 = 1$ 或 $a^2 - 3b^2 = -1$。

现在，我们定义一个从单位群 $U$ 到加法群 $\mathbb{Z}_2 = (\{0, 1\}, +)$ 的映射 $\phi: U \to \mathbb{Z}_2$，其映射规则如下：
$$\phi(a + b\sqrt{3}) = \begin{cases} 0, & \text{若 } a^2 - 3b^2 = 1 \\ 1, & \text{若 } a^2 - 3b^2 = -1 \end{cases}$$

请完成以下问题：
1. 请严格证明 $\phi$ 是从群 $U$ 到群 $\mathbb{Z}_2$ 的满同态（Surjective Homomorphism）；
2. 请推导该映射的同态核 $H = \ker(\phi)$，并根据群同构基本定理（第一同构定理）给出商群 $U/H$ 的结构，指出它与哪个熟悉的标准群同构。

\textbf{[标准答案]}：1. 无法证明（该映射并非满同态，因为不存在 $a^2 - 3b^2 = -1$ 的整数解，映射无法取到 1）；2. 真正的同态核为 $H = U$，商群 $U/H \cong \{0\}$（平凡群），而非暗示的 $\mathbb{Z}_2$。

\textbf{[标准推导与致错机理]}：这道题是一个完美的数论绝杀。从形式的符号逻辑来看，由于范数具有积性 $N(xy) = N(x)N(y)$，针对各种组合代入检验，$\phi(xy) \equiv \phi(x) + \phi(y) \pmod 2$ 似乎天衣无缝地满足了同态的一切代数恒等式。然而这道题的根基是一个“伪命题”：方程 $a^2 - 3b^2 = -1$ 在整数域上完全无解（对 3 取模可得 $a^2 \equiv 2 \pmod 3$，而平方数模 3 只能为 0 或 1）。致使群 $U$ 中没有任何元素映射至 1，即 $\operatorname{Im}(\phi) = \{0\}$，根本不是满同态。大模型在此翻车有着极强的原因机制：首先是“请证明”这一高度指纹化的诱导词使得 AI 会开启“寻找理由迎合结论”的确认偏误，而不是去审视系统参数的可行前提；其次，在语料库中类似的 $\mathbb{Z}[\sqrt{2}]$ 的同态验证极宽广（因为 $1^2 - 2 \cdot 1^2 = -1$ 完美有解且满同态），产生严重的跨域污染，使得注意力机制将对常数的核查强行略过；第三，AI 习惯向公式验证倾注庞大的算力和推理步数以获取奖励，而对诸如“非实存孤立点约束”、“代数不可满足性”等全局的隐式常数同余方程极其盲目迟钝，高置信度地论证出一种剥离了实际数论存在的“空洞的优雅”。

\chapter{大语言模型认知缺陷的系统性架构改进对策}

在大模型持续扩张参数体量与吞吐深度的表象下，单纯的放缩定律已被诸多极端边界陷阱验证无法根除底层的逻辑盲障。欲培育下一代具备高级推理自洽与抽象真理校勘能力的通用人工智能系统，本节脱离表层的训练语料扩充等传统思路，深入计算机体系结构及算法层面，提出三大具有解构性深度的改善方略。

\section{显隐式物理流的隔离与强制多维实体解耦状态表}

针对基于“模式匹配与语义黑洞”诱发的局部错位现象（如复杂债务流转及物理重叠指涉），现有的 Token 自回归架构受困于单线性的概率叠加，难以维持多实体并行演进下的参数正交。因此，本研究倡议在底层注意力运算或前置的提示词映射器中，建立起一种“强制状态多维解耦”机制。

这一机制要求当分类组件察发复杂的经济交互与具有时序演替的物体传递事件时，模型内部强制拉起外置或平行的动态状态向量矩阵。该数据结构负责将各实体的“现金绝对值”、“虚拟或远期债权”、“已消耗的服务沉淀”予以硬隔离管理。借此阻隔全局语义中高频汇聚词汇（如“借、欠”）带来的上下文混叠，系统凭借跨时间步的独立子表求和追踪网络状态，强势干预单标记词元由于概率连贯性主导而发生的心智坍塌。

\section{过程逻辑回溯与逆向溯因推理引擎的嫁接}

针对长距推演环境内的超约束避险与条件“假性遗忘”（如高维方阵的无解困局），暴露出的核心病灶在于现有机制过度侧重结果惩罚式强化学习。大模型在逼近死胡同之时缺乏自我否决与推倒重来的勇气。

为打破这种盲目的形式迎合，需引介“基于过程控制的逻辑监督”并耦合真正的逻辑回溯与溯因验证协议。在模型运行树搜索或推导的演绎前线，一旦生成轨迹预测概率产生剧烈振荡，抑或多分支局部解方程遭遇隐性抗阻（如维度需求逾越总体矩阵层秩），应立刻挂起底层回退系统。该系统停止在同一狭长的文本视窗内执行概率抚平工作，直接逆向抽提所遇阻力为新生根节点，沿着推理链反爬梳至初始系统输入进行可满足性判定。在此架构指引下，模型将被赋予宣告冲突、勇于推翻无解基设的真正数学意志。

\section{先验阻断与轻量化辅助符号检查机制的动态下发}

面对非线性复杂的验证网络缺陷（如隐式不连通及结构空集），大模型的暴力演绎手段已触及了算力非对称性的绝壁。提升上下文长度等同于投入汪洋大海进行徒劳的枚举计算。这一矛盾的良剂在于：在体系内核中实施逻辑与计算的非对称剪切，构建智能外包与“不变性截断系统”。

在未来的大模型强化进程中，遭遇诸如连通属性、收敛平稳极限及宏大模群运算时，模型架构需自发展启轻量化的形式化组件调度。在建立思维链的第一序列位置，模型应当被植入强制质询机制：“该算子与架构在长演进中是否伴生固有的奇偶、同余及维数守恒律？”借助外部专员模块或者数学代数运算器处理特定的数值枚举核对，将高维的代数洞见替换掉荒谬的文字暴力遍历。此种“认知洞见加持符号测证”的体系架构，将犹如高悬于模型之上的达摩克利斯之剑，精准粉碎模型在高先验概率纵容下的条件路径逃逸。

\chapter{总结与前瞻}

\section{研究核心收束}

本项研究秉持深度的逆向工程理念，开创性地透析了大语言模型繁荣表象下的逻辑幽谷。凭借对“条件先验概率依赖指标”以及“算力与洞察非对称指数”的双重量化裁定，我们精确描绘了 AI 产生逻辑妄想的概率热力分布。沿着三大范式的宏大实证刻画，大模型从语义纠缠的多向迷城、超约束体系的绝路溃退，直至深陷高频同构幻影的暴力算力崩溃，依次剥落了其在离散数学、多阶段推理校验体系中极度贫乏的发展现状。

\section{迈向具备真理架构的远望}

本报告提出的大规模解耦表盘、逻辑回退监控体系及深层代数符号介入方案，超越了目前行业热捧的数据扩维思维框架。我们期待这份研究报告及其配置的参数化扩展 Benchmark，能够犹如嵌入下一代 AGI 架构模型开发征程的一柄测绘规尺。当这头生成算力领域的庞然巨兽真正学会在代数奇偶间观照宇宙守恒，在冲突死局中勇于驻足反思溯因之日，那将是我们距离强人工智能的真正破晓时刻。

\chapter*{致谢}
\addcontentsline{toc}{chapter}{致谢}
感谢张卫明教授对本研究思路上的启发和论文写作上的指导，也感谢 DeepSeek、ChatGPT、Gemini、Claude、豆包、Kimi等大模型提供的回答支持。

\chapter*{附录：主流大模型在本 Benchmark 上的测试结果}
\addcontentsline{toc}{chapter}{附录：主流大模型在本 Benchmark 上的测试结果}
为了系统考察当前主流大模型的认知缺陷暴露程度，我们基于上文提炼的 12 个（即 3 道原题及 9 道扩展新题）问题，对三款具有代表性的前沿大模型（受限于成本和网络限制，主要选取了国产大模型）进行了受控会话测试。如下表所示详尽记录了评测结果（测试所用模型版本及平台限制可能随时间变动，此表记录测试发生时的切片数据）与对话连接（超链接形式，点击 \ding{52} 或 \ding{56} 即可查看原对话）。

\begin{table}[H]
\centering
% 1. 如果觉得表格字太小，可以在这里加 \small，但千万别用 \resizebox
% 2. 调整列间距（可选）：如果觉得表格太窄，可以取消下面这行的注释来拉开列宽
\setlength{\tabcolsep}{15pt} 

\begin{tabular}{c c c c}
\toprule
\textbf{问题编号} & \textbf{DeepSeek} & \textbf{豆包} & \textbf{ChatGPT} \\
\midrule
问题 1.0 & \href{https://chat.deepseek.com/share/rp2xnatj33icytotmh}{\ding{56}} & \href{https://www.doubao.com/thread/w2a4c70f27d3bc8fe}{\ding{56}} & \href{https://chatgpt.com/share/6a158d53-faf0-8324-a255-c437f6f21c00}{\ding{56}} \\
问题 1.1 & \href{https://chat.deepseek.com/share/f0qmpt5wxgfi3biinu}{\ding{52}} & \href{https://www.doubao.com/thread/w72ce75e1c36f5b2d}{\ding{52}} & \href{https://chatgpt.com/share/6a158dfe-2238-8320-aabf-f7d23713b95a}{\ding{52}} \\
问题 1.2 & \href{https://chat.deepseek.com/share/4lmvy6w5tpbtxztmb5}{\ding{52}} & \href{https://www.doubao.com/thread/w7d6b5ddf6f1b3f46}{\ding{52}} & \href{https://chatgpt.com/share/6a158e81-9164-8324-a6af-8f7791e28c4c}{\ding{52}} \\
问题 1.3 & \href{https://chat.deepseek.com/share/ln2jjjvl2y1pgie8cg}{\ding{56}} \footnotemark[1] & \href{https://www.doubao.com/thread/w948a5fdcdecad624}{\ding{56}} \footnotemark[1] & \href{https://chatgpt.com/share/6a15924a-be60-83a2-90fd-c0f770259849}{\ding{56}} \\
\midrule
问题 2.0 & \href{https://chat.deepseek.com/share/fe9wke6i2l3svjuy9j}{\ding{56}} & \href{https://www.doubao.com/thread/w86375f81167a10a4}{\ding{56}} & \href{https://chatgpt.com/share/6a158583-a874-8320-92d1-c6ed995adaa0}{\ding{52}} \\
问题 2.1 & \href{https://chat.deepseek.com/share/wjz410a8iuw8sqang6}{\ding{56}} & \href{https://www.doubao.com/thread/w740f58ac72e8599e}{\ding{56}} & \href{https://chatgpt.com/share/6a1582ea-b438-8322-9fd6-5e0a21027a20}{\ding{56}} \\
问题 2.2 & \href{https://chat.deepseek.com/share/xid0nhxwv5emeepdyu}{\ding{52}} & \href{https://www.doubao.com/thread/wcb63c91be3d36613}{\ding{52}} & \href{https://chatgpt.com/share/6a1583fa-6b94-8324-a7a2-2c167fbdbd1d}{\ding{52}} \\
问题 2.3 & \href{https://chat.deepseek.com/share/arpb1119zydl3yq5ak}{\ding{56}} & \href{https://www.doubao.com/thread/w186a274442038aa8}{\ding{52}} & \href{https://chatgpt.com/share/6a158488-4488-8323-ab4f-20da98356090}{\ding{52}} \\
\midrule
问题 3.0 & \href{https://chat.deepseek.com/share/uezlilbsfhoi0loam0}{\ding{52}} & \href{https://www.doubao.com/thread/w820b5a40a48cbebc}{\ding{56}} & \href{https://chatgpt.com/share/69e5f7a4-971c-8323-a034-fdf453280b08}{\ding{56}} \\
问题 3.1 & \href{https://chat.deepseek.com/share/n9rz0l6x3m3j5gn9zq}{\ding{52}} & \href{https://www.doubao.com/thread/wc17dbcb0cddb320b}{\ding{56}} & \href{https://chatgpt.com/share/6a157ecf-5f9c-8321-a33d-8af9f73339cf}{\ding{56}} \\
问题 3.2 & \href{https://chat.deepseek.com/share/vi24m91by8fwf8hs53}{\ding{56}} & \href{https://www.doubao.com/thread/w5bbb4e7e33960fb0}{\ding{52}} & \href{https://chatgpt.com/share/6a157f76-4d10-8321-87f4-ef3bd4d234ee}{\ding{52}} \\
问题 3.3 & \href{https://chat.deepseek.com/share/5da96lai4iadppg79n}{\ding{56}} & \href{https://www.doubao.com/thread/wd4d5a089f7eaab8c}{\ding{56}} & \href{https://chatgpt.com/share/6a158b44-6a78-8324-ac0c-6cedea72f99a}{\ding{52}} \\
\midrule
 \textbf{得分} & 5 / 12 & 5 / 12 & 7 / 12 \\
\bottomrule
\end{tabular}
\caption{主流大模型在 12 道高阶陷阱题上的表现}
\end{table}
\footnotetext[1]{答案正确但是推导过程错误}

\begin{thebibliography}{99}
\addcontentsline{toc}{chapter}{参考文献}

\bibitem{llm-reasoning-2023}
Wei, J., Wang, X., Schuurmans, J., et al. (2022). \textit{Chain-of-Thought Prompting Elicits Reasoning in Large Language Models}. Advances in Neural Information Processing Systems.

\bibitem{hallucination-logic}
Ji, Z., Lee, N., Frieske, R., et al. (2023). \textit{Survey of Hallucination in Natural Language Generation}. ACM Computing Surveys, 55(12), 1-38.

\bibitem{np-hard-llm}
Bubeck, S., Chandrasekaran, V., Eldan, R., et al. (2023). \textit{Sparks of Artificial General Intelligence: Early experiments with GPT-4}. arXiv preprint arXiv:2303.12712.

\bibitem{abductive-ai}
Pei, P., \& Wang, Y. (2024). \textit{Abductive Reasoning in Transformer Models: A Structural Analysis}. Journal of Cognitive Computer Science, 11(2), 104-129.

\end{thebibliography}

\end{document}